\documentclass[
10pt, % Main document font size
a4paper, % Paper type, use 'letterpaper' for US Letter paper
oneside, % One page layout (no page indentation)
%twoside, % Two page layout (page indentation for binding and different headers)
headinclude,footinclude, % Extra spacing for the header and footer
BCOR5mm, % Binding correction
]{scrartcl}
%%%%%%%%%%%%%%%%%%%%%%%%%%%%%%%%%%%%%%%%%
% Arsclassica Article
% Structure Specification File
%
% This file has been downloaded from:
% http://www.LaTeXTemplates.com
%
% Original author:
% Lorenzo Pantieri (http://www.lorenzopantieri.net) with extensive modifications by:
% Vel (vel@latextemplates.com)
%
% License:
% CC BY-NC-SA 3.0 (http://creativecommons.org/licenses/by-nc-sa/3.0/)
%
%%%%%%%%%%%%%%%%%%%%%%%%%%%%%%%%%%%%%%%%%

%----------------------------------------------------------------------------------------
%	REQUIRED PACKAGES
%----------------------------------------------------------------------------------------

\usepackage[
nochapters, % Turn off chapters since this is an article        
beramono, % Use the Bera Mono font for monospaced text (\texttt)
eulermath,% Use the Euler font for mathematics
pdfspacing, % Makes use of pdftex’ letter spacing capabilities via the microtype package
dottedtoc % Dotted lines leading to the page numbers in the table of contents
]{classicthesis} % The layout is based on the Classic Thesis style

\usepackage{arsclassica} % Modifies the Classic Thesis package
\usepackage{ulem}
\usepackage{graphicx}
\usepackage{float}
\usepackage{algorithm2e}
\usepackage[export]{adjustbox}
\usepackage[utf8]{inputenc}  
\usepackage[english,frenchb]{babel} 
\usepackage[T1]{fontenc} % Required for accented characters
\usepackage{multicol}
\usepackage{graphicx} % Required for including images
\graphicspath{{Figures/}} % Set the default folder for images

\usepackage{enumitem} % Required for manipulating the whitespace between and within lists

\usepackage{lipsum} % Used for inserting dummy 'Lorem ipsum' text into the template


\expandafter\def\csname ver@subfig.sty\endcsname{}


\usepackage{svg}
\usepackage{hyperref}

\usepackage{amsmath,amssymb,amsthm} % For including math equations, theorems, symbols, etc

\usepackage{varioref} % More descriptive referencing

%----------------------------------------------------------------------------------------
%	THEOREM STYLES
%---------------------------------------------------------------------------------------

\theoremstyle{definition} % Define theorem styles here based on the definition style (used for definitions and examples)
\newtheorem{definition}{Definition}

\theoremstyle{plain} % Define theorem styles here based on the plain style (used for theorems, lemmas, propositions)
\newtheorem*{remark}{Remark}
\newtheorem{theorem}{Theorème}
\newtheorem{lemma}[theorem]{Lemme}
\newtheorem{proposition}[theorem]{Proposition}

\theoremstyle{remark} % Define theorem styles here based on the remark style (used for remarks and notes)

%----------------------------------------------------------------------------------------
%	HYPERLINKS
%---------------------------------------------------------------------------------------

\hypersetup{
%draft, % Uncomment to remove all links (useful for printing in black and white)
colorlinks=true, breaklinks=true, bookmarks=true,bookmarksnumbered,
urlcolor=webbrown, linkcolor=RoyalBlue, citecolor=webgreen, % Link colors
pdftitle={}, % PDF title
pdfauthor={\textcopyright}, % PDF Author
pdfsubject={}, % PDF Subject
pdfkeywords={}, % PDF Keywords
pdfcreator={pdfLaTeX}, % PDF Creator
pdfproducer={LaTeX with hyperref and ClassicThesis} % PDF producer
}

\hyphenation{Fortran hy-phen-ation} 

\title{\normalfont{FTML Exercices 5 solutions}\\
} 
\author{\spacedlowsmallcaps{}
} 
\date{\today}
\begin{document}
\renewcommand{\sectionmark}[1]{\markright{\spacedlowsmallcaps{#1}}}
\lehead{\mbox{\llap{\small\thepage\kern1em\color{halfgray} \vline}\color{halfgray}\hspace{0.5em}\rightmark\hfil}} 
\pagestyle{scrheadings}
\maketitle 
\setcounter{tocdepth}{3} 

\tableofcontents

\section{\large\color{Blue}Logistic regression}

1]

\begin{equation}
    \begin{aligned}
        \label{eq:}
	\forall z\in \mathbb{R}, \sigma'(z) &= \big(- \frac{1}{(1+e^{-z})^2}\big)\big( -e^{-z}\big)\\
	&= \frac{1}{1+e^{-z}} \frac{e^{-z}}{1+e^{-z}} \\
	&= \frac{1}{1+e^{-z}} \frac{e^{-z}e^z}{(1+e^{-z})e^z} \\
	&= \frac{1}{1+e^{-z}} \frac{1}{1+e^{z}} \\
	&= \sigma(z)\sigma(-z)
    \end{aligned}
\end{equation}


2] We compute the second order derivative.

\begin{equation}
    \begin{aligned}
        \label{eq:}
	\frac{\partial l}{\partial \hat{y}}( \hat{y}, y) &= \frac{-ye^{- \hat{y}y}}{1+e^{- \hat{y}y}}\\
	&= \frac{-ye^{- \hat{y}y}}{1+e^{- \hat{y}y}}\frac{e^{\hat{y}y}}{e^{\hat{y}y}}\\
	&= \frac{-y}{e^{\hat{y}y}+1}\\
	&= -y\sigma(-\hat{y}y)
    \end{aligned}
\end{equation}

Hence,

\begin{equation}
    \begin{aligned}
        \label{eq:}
	\frac{\partial^2 l}{\partial \hat{y}^2}( \hat{y}, y) &= -y\sigma(-\hat{y}y)\sigma(\hat{y}y)\times -y\\
	&= y^2\sigma(-\hat{y}y)\sigma(\hat{y}y)>0\\
    \end{aligned}
\end{equation}


Hence, the second-order derivative is strictly positive, and $l( \hat{y}, y)$ is stricly convex in its first argument.
\\

3] We introduce the following functions:

$$
g_i = \left\{
    \begin{array}{ll}
	\mathbb{R}^d\rightarrow \mathbb{R}    \\
	\theta \mapsto   l(x_i^T\theta, y_i) \\
    \end{array}
    \right.
    $$

$$
u_i = \left\{
    \begin{array}{ll}
	\mathbb{R}\rightarrow \mathbb{R}    \\
	\hat{y} \mapsto   l(\hat{y}, y_i) \\
    \end{array}
    \right.
    $$

$$
v_i = \left\{
    \begin{array}{ll}
	\mathbb{R}^d\rightarrow \mathbb{R}    \\
	\theta \mapsto   x_i^T\theta \\
    \end{array}
    \right.
    $$

    Then, $\forall i$

    \begin{equation}
	\label{eq:comp}
	l(x_i^T\theta, y_i) = g_i(\theta)= (u_i\circ v_i)(\theta)
    \end{equation}

\textbf{a]} (convexity) It is sufficient to show that each $g_i:\theta \rightarrow l(x_i^T\theta, y_i) $ is convex, because the sum of convex functions is convex. By definition, (equation~\ref{eq:comp}), $g_i$ is a convex function $u_i$ applied to a linear mapping $v_i$, which prooves that $g_i$ is convex.
\\

\textbf{b]} (gradient) By composition of the jacobian matrices,
    \begin{equation}
    L_{\theta}^{g_i}= L_{v_i(\theta)}^{u_i}L_{\theta}^{v_i}= u_i'(v_i(\theta))L_{\theta}^{v_i}
    \end{equation}

    Or equivalently:
    \begin{equation}
	\nabla_{\theta}g_i(\theta) = u_i'(v_i(\theta))\nabla_{\theta}v_i(\theta)
    \end{equation}

    We already know that $\nabla_{\theta}v_i(\theta) = x_i$.
    \\

	 In question 2, we have seen that $\forall y, \hat{y}$,
	    \begin{equation}
		\frac{\partial l}{\partial \hat{y}}( \hat{y}, y) = -y\sigma(-\hat{y}y)
	    \end{equation}

	    Hence,
	    \begin{equation}
		\begin{aligned}
		    u_i'(v_i(\theta))&= -y_i\sigma( -v_i(\theta)y_i)\\
		    &= -y_i\sigma( -x_i^T\theta y_i)\\
		\end{aligned}
	    \end{equation}

	Finally,

    \begin{equation}
	\nabla_{\theta}g_i(\theta) = -y_i\sigma( -x_i^T\theta y_i)x_i
    \end{equation}
    

    and 

    \begin{equation}
	\begin{aligned}
	    \label{eq:}
	    \nabla_{\theta}R_n(\theta) &= \frac{1}{n} \sum^{n}_{i=1}\nabla_{\theta}g_i(\theta) \\
        &= \frac{1}{n}  \sum^{n}_{i=1} -y_i\sigma(-x_i^T\theta y_i)x_i
	\end{aligned}
    \end{equation}


\end{document}
